% GENERATED FILE. Edit canonical lessons, then run npm run generate:books.
% canonical-source-hash: fnv1a64:2cc36e8043b81011
% canonical-lessons: LA-C341-obsideo, LA-C341-obsidio, LA-C341-oppugno, LA-C341-expugno, LA-C341-aries

\chapter{Obsidiō}
\label{ch:la-obsidio}
\hlchaptermodality{341}

\begin{chapteropening}
\textbf{By the end of this chapter:} \emph{I can tell how a town was besieged: to besiege, the siege, to attack, to take by storm and the battering ram.}

Complete the last lesson of chapter 341: I can tell how a town was besieged: to besiege, the siege, to attack, to take by storm and the battering ram.
\end{chapteropening}

\section[obsideō, obsidēre]{obsideō, obsidēre — to besiege; to sit at}

\label{lesson:LA-C341-obsideo}



\begin{quote}
Before the new one: say the Latin for a breastplate, then the Latin for a round shield.
\end{quote}

\subsection*{You'll want to know: obsideō, obsidēre}
\textbf{obsideō, obsidēre} — \textquotedblleft{}to besiege; to sit at\textquotedblright{}.

\begin{quote}
Hostēs oppidum trēs mēnsēs obsēdērunt.
\end{quote}

\begin{quote}
\emph{The enemy besieged the town for three months.}
\end{quote}

\emph{obsideō, obsidēre, obsēdī, obsessum}: \emph{ob} + \emph{sedeō} (taught), to sit down in front of. English \emph{obsess}.

\begin{cousinweb}
\textbf{obsideō} is built on \emph{sedeō}, \textquotedblleft{}to sit\textquotedblright{}, with \emph{ob}, \textquotedblleft{}in front of\textquotedblright{} (Lewis \& Short): to sit down before a town. English \emph{obsess} comes from its participle.
\end{cousinweb}

\begin{grammarlens}[title={What you've built}]
One more word for telling, explaining and describing at B1.
\end{grammarlens}

\subsection*{Your turn}
\begin{itemize}
\raggedright
  \item \emph{Say it:} \emph{obsideō, obsidēre}
  \item \emph{Say it:} \emph{obsideō, obsidēre}, once more
  \item \emph{Say it:} a sentence of your own with \emph{obsideō, obsidēre}
  \item \emph{Recall:} say the Latin for a breastplate, then the Latin for a round shield, then say \emph{obsideō, obsidēre} again
\end{itemize}

\begin{tcolorbox}[breakable,colback=teal!4,colframe=teal!35!black,title={Before you move on}]
What does obsideō, obsidēre mean? (\textquotedblleft{}To besiege; to sit at\textquotedblright{}.) Say it once more.
\end{tcolorbox}
\section[obsidiō, obsidiōnis]{obsidiō, obsidiōnis — a siege}

\label{lesson:LA-C341-obsidio}



\begin{quote}
Before the new one: say the Latin for a round shield, then the Latin for to besiege.
\end{quote}

\subsection*{You'll want to know: obsidiō, obsidiōnis}
\textbf{obsidiō, obsidiōnis} — \textquotedblleft{}a siege\textquotedblright{}.

\begin{quote}
Obsidiō longa erat, et cīvēs famē labōrābant.
\end{quote}

\begin{quote}
\emph{The siege was long, and the citizens were suffering from hunger.}
\end{quote}

Feminine, 3rd declension, from \emph{obsideō}.

\begin{cousinweb}
\textbf{obsidiō} is built on \emph{obsideō}, \textquotedblleft{}to besiege\textquotedblright{} (EtymDB), with the ending \emph{-iō}.
\end{cousinweb}

\begin{grammarlens}[title={What you've built}]
One more word for telling, explaining and describing at B1.
\end{grammarlens}

\subsection*{Your turn}
\begin{itemize}
\raggedright
  \item \emph{Say it:} \emph{obsidiō, obsidiōnis}
  \item \emph{Say it:} \emph{obsidiō, obsidiōnis}, once more
  \item \emph{Say it:} a sentence of your own with \emph{obsidiō, obsidiōnis}
  \item \emph{Recall:} say the Latin for a round shield, then the Latin for to besiege, then say \emph{obsidiō, obsidiōnis} again
\end{itemize}

\begin{tcolorbox}[breakable,colback=teal!4,colframe=teal!35!black,title={Before you move on}]
What does obsidiō, obsidiōnis mean? (\textquotedblleft{}A siege\textquotedblright{}.) Say it once more.
\end{tcolorbox}
\section[oppugnō, oppugnāre]{oppugnō, oppugnāre — to attack, to assault}

\label{lesson:LA-C341-oppugno}



\begin{quote}
Before the new one: say the Latin for to besiege, then the Latin for a siege.
\end{quote}

\subsection*{You'll want to know: oppugnō, oppugnāre}
\textbf{oppugnō, oppugnāre} — \textquotedblleft{}to attack, to assault\textquotedblright{}.

\begin{quote}
Mīlitēs mūrōs oppidī oppugnābant.
\end{quote}

\begin{quote}
\emph{The soldiers were attacking the walls of the town.}
\end{quote}

\emph{oppugnō, oppugnāre, oppugnāvī, oppugnātum}: \emph{ob} + \emph{pugnō} (taught). To attack, not yet to win.

\begin{cousinweb}
\textbf{oppugnō} is built on \emph{pugnō}, \textquotedblleft{}to fight\textquotedblright{}, with \emph{ob-}, \textquotedblleft{}against\textquotedblright{} (EtymDB). English \emph{oppugn} comes from it.
\end{cousinweb}

\begin{grammarlens}[title={What you've built}]
One more word for telling, explaining and describing at B1.
\end{grammarlens}

\subsection*{Your turn}
\begin{itemize}
\raggedright
  \item \emph{Say it:} \emph{oppugnō, oppugnāre}
  \item \emph{Say it:} \emph{oppugnō, oppugnāre}, once more
  \item \emph{Say it:} a sentence of your own with \emph{oppugnō, oppugnāre}
  \item \emph{Recall:} say the Latin for to besiege, then the Latin for a siege, then say \emph{oppugnō, oppugnāre} again
\end{itemize}

\begin{tcolorbox}[breakable,colback=teal!4,colframe=teal!35!black,title={Before you move on}]
What does oppugnō, oppugnāre mean? (\textquotedblleft{}To attack, to assault\textquotedblright{}.) Say it once more.
\end{tcolorbox}
\section[expugnō, expugnāre]{expugnō, expugnāre — to take by storm}

\label{lesson:LA-C341-expugno}



\begin{quote}
Before the new one: say the Latin for a siege, then the Latin for to attack.
\end{quote}

\subsection*{You'll want to know: expugnō, expugnāre}
\textbf{expugnō, expugnāre} — \textquotedblleft{}to take by storm\textquotedblright{}.

\begin{quote}
Tandem oppidum expugnāvērunt.
\end{quote}

\begin{quote}
\emph{At last they took the town by storm.}
\end{quote}

\emph{expugnō, expugnāre, expugnāvī, expugnātum}. After the attack comes, if things go well, the storming.

\begin{cousinweb}
\textbf{expugnō} is built on \emph{pugnō}, \textquotedblleft{}to fight\textquotedblright{}, with \emph{ex-}, \textquotedblleft{}out, to the end\textquotedblright{} (Lewis \& Short): to fight a place out of the enemy's hands.
\end{cousinweb}

\begin{grammarlens}[title={What you've built}]
One more word for telling, explaining and describing at B1.
\end{grammarlens}

\subsection*{Your turn}
\begin{itemize}
\raggedright
  \item \emph{Say it:} \emph{expugnō, expugnāre}
  \item \emph{Say it:} \emph{expugnō, expugnāre}, once more
  \item \emph{Say it:} a sentence of your own with \emph{expugnō, expugnāre}
  \item \emph{Recall:} say the Latin for a siege, then the Latin for to attack, then say \emph{expugnō, expugnāre} again
\end{itemize}

\begin{tcolorbox}[breakable,colback=teal!4,colframe=teal!35!black,title={Before you move on}]
What does expugnō, expugnāre mean? (\textquotedblleft{}To take by storm\textquotedblright{}.) Say it once more.
\end{tcolorbox}
\section[ariēs, arietis]{ariēs, arietis — a battering ram; a ram}

\label{lesson:LA-C341-aries}



\begin{quote}
Before the new one: say the Latin for to attack, then the Latin for to take by storm.
\end{quote}

\subsection*{You'll want to know: ariēs, arietis}
\textbf{ariēs, arietis} — \textquotedblleft{}a battering ram; a ram\textquotedblright{}.

\begin{quote}
Ariete portam frēgērunt.
\end{quote}

\begin{quote}
\emph{They broke the gate with a battering ram.}
\end{quote}

Masculine, 3rd declension. First the male sheep, then the beam with an iron ram's head that broke gates.

\begin{cousinweb}
\textbf{ariēs} is inherited; EtymDB traces it to a Proto-Indo-European name for a domestic animal, with cousins in Greek \emph{eriphos}, \textquotedblleft{}kid\textquotedblright{}.
\end{cousinweb}

\begin{grammarlens}[title={What you've built}]
One more word for telling, explaining and describing at B1.
\end{grammarlens}

\subsection*{Your turn}
\begin{itemize}
\raggedright
  \item \emph{Say it:} \emph{ariēs, arietis}
  \item \emph{Say it:} \emph{ariēs, arietis}, once more
  \item \emph{Say it:} a sentence of your own with \emph{ariēs, arietis}
  \item \emph{Recall:} say the Latin for to attack, then the Latin for to take by storm, then say \emph{ariēs, arietis} again
\end{itemize}

\begin{tcolorbox}[breakable,colback=teal!4,colframe=teal!35!black,title={Before you move on}]
What does ariēs, arietis mean? (\textquotedblleft{}A battering ram; a ram\textquotedblright{}.) Say it once more.
\end{tcolorbox}
